\section{Discussion and limitations}\label{sec:discussion}

\paragraph{Evaluation posture.} The closed-loop results are in silico.
The plant is synthetic, though every physical ingredient, from
solubility and kinetics to the brine envelope and the actuator family,
is grounded in primary literature and, wherever independent data
exists, validated out-of-sample component by component
(Section~\ref{sec:validation}). We regard component-wise real-data
validation combined with a parameterized-mismatch closed loop as the
honest intermediate between simulation-only studies and a field trial.
A field pilot is the natural next step, and it would require exactly
the sparse instrumentation this architecture is designed around.

\paragraph{Constraint units.} The scaling constraint is expressed as a
polymerization (monomer-depletion) rate rather than a wall-deposition
thickness rate. The molecular-deposition treatment of
\citet{weres1981kinetics}, transcribed under the same integrity
protocol (its printed checkpoints reproduce, including the 3.0~kcal/mol
effective activation energy; one internal inconsistency in its
Eq.~[4c] is documented rather than silently repaired), closes the unit
gap. Along the envelope where the constraint is active,
$r_{\max} = 0.5$~g\,t$^{-1}$\,min$^{-1}$ corresponds to a Weres-implied
molecular-deposition wall growth, a lower bound on total scale growth,
of about 0.3~mm/yr at the optimized operating point and 0.04 to
1.6~mm/yr across 90--175~$^{\circ}$C; the four-fold $r_{\max}$ band of
the sensitivity study spans that calibration.

\paragraph{Scenario fencing.} Trade-press figures inform scenario scale
only. The rule is enforced mechanically in the repository, and every
peso and megawatt headline carries an envelope range with decision
invariance verified rather than assumed.

\paragraph{Design constants.} The back-off gain is fixed at
$\kappa = 1$. Its ablation over $\kappa \in \{0,1\}$ brackets the
safety--harvest trade under stress, and $\sigma_{\mathrm{LOO}}$ is
already data-adaptive. A continuous $\kappa$ sweep is the tuning study we
would prioritize next.

\paragraph{Scope boundaries.} High-pH brines require silica ionization
corrections that are deferred here; the speciation layer exists in the
repository but is not exercised in the loop. The identifiability
principle of Section~\ref{sec:sensorval} is likewise a scope statement
rather than a weakness: corrections are used only where labels can
support them, and the architecture knows which regime it is in.

\paragraph{What generalizes.} Nothing in
Eqs.~\eqref{eq:sensor}--\eqref{eq:rto} is specific to silica. The
pattern, a physics prior corrected only to its earned credibility and
feeding a constraint backed off by the correction's own cross-validated
uncertainty, applies to any hard-constrained process operated against a
structurally wrong model: fouling, crystallization, corrosion, membrane
scaling. Geothermal silica supplies an unusually clean and unusually
demanding instance, with quartic constraint sensitivity, an induction
delay, non-monotonic actuation, and a measured sixfold model-family
gap, and that is why we offer it as a benchmark.
